\documentclass[11pt]{article}
\usepackage{nmrdoc}
\usepackage{listings}

\lstset{
  basicstyle=\ttfamily\small,
  columns=fullflexible,
  keepspaces=true,
  breaklines=true,
  breakatwhitespace=true,
  frame=single,
  xleftmargin=0pt,
  xrightmargin=0pt
}
\setlength{\parindent}{0pt}
\setlength{\parskip}{0.65em}
\emergencystretch=3em
\AtBeginDocument{\raggedright}

\begin{document}
\NmrTitle{\texttt{nmr\_extract} Operator Use Guide}

\section{What the extractor does}

\texttt{nmr\_extract} builds a charged protein model from one of five input
shapes, runs the calculator stack for that model, and writes NumPy feature
arrays. In trajectory mode it also streams GROMACS frames, writes per-frame
outputs, and writes a trajectory HDF5 file.

APBS and AIMNet2 are part of the run in every mode. They are not controlled by
\texttt{--no-apbs}, \texttt{--no-coulomb}, or any other switch. MOPAC is the
available expensive-method toggle, but the current parser implements it
asymmetrically: modes 1--4 default MOPAC on and accept
\texttt{--no-mopac}; mode 5 defaults MOPAC off and accepts \texttt{--mopac}.

\section{Command shape}

Every run needs exactly one mode selector:

\begin{lstlisting}
--pdb FILE
--protonated-pdb FILE
--orca --root NAME
--mutant --wt NAME --ala NAME
--trajectory DIR
\end{lstlisting}

The parser rejects a command with no mode:

\begin{lstlisting}
no mode flag found
  valid modes: --pdb, --protonated-pdb, --orca, --mutant, --trajectory
\end{lstlisting}

It also rejects more than one mode selector:

\begin{lstlisting}
multiple mode flags present: MODE1 and MODE2 (pick exactly one)
\end{lstlisting}

\texttt{--help}, \texttt{-h}, or no arguments prints usage and exits without
running. File existence is checked after parsing. For a required file, the path
must exist and must be a regular file. For a required directory, the path must
exist and must be a directory.

Unknown flags are not a supported interface. The current parser does not have
a general unknown-flag rejection pass, so an unknown flag can be ignored when a
valid mode is otherwise present. Do not rely on that behavior.

\section{Common options}

\texttt{--output DIR} is required for every successful run. The directory may
already exist. If it exists, it must be a directory:

\begin{lstlisting}
--output: exists but is not a directory: DIR
\end{lstlisting}

If \texttt{--output} is omitted, execution stops with:

\begin{lstlisting}
ERROR: --output DIR required
\end{lstlisting}

\texttt{--config FILE} is optional. It points to a TOML file with calculator
parameter overrides. If provided, validation requires it to be a regular file:

\begin{lstlisting}
--config: file not found: FILE
--config: not a regular file: FILE
\end{lstlisting}

If \texttt{--config} is not provided, the executable tries
\texttt{NMR\_DATA\_DIR/calculator\_params.toml} when that file exists.

\texttt{--aimnet2 FILE} is an AIMNet2 model path, not an AIMNet2 on/off
switch. AIMNet2-derived calculators require a model. The resolution order is:
the CLI path, then \texttt{NMR\_AIMNET2\_MODEL}, then
\texttt{aimnet2\_model\_path} from the loaded calculator TOML. A TOML-relative
model path is resolved relative to that TOML file. If a CLI or resolved path
is present, validation requires a regular file:

\begin{lstlisting}
--aimnet2 model: file not found: FILE
--aimnet2 model: not a regular file: FILE
\end{lstlisting}

If no model path can be resolved, execution stops before any mode runner:

\begin{lstlisting}
ERROR: AIMNet2 model required: pass --aimnet2 PATH, set NMR_AIMNET2_MODEL, or set aimnet2_model_path in calculator_params.toml.
\end{lstlisting}

\section{Mode 1: bare PDB}

Use this mode for a PDB that should be protonated by \texttt{reduce} at runtime.

\begin{lstlisting}
nmr_extract --pdb input.pdb --output out --aimnet2 model.jpt
\end{lstlisting}

Optional flags:

\begin{itemize}
\item \texttt{--pH VALUE}: parsed as a floating-point value with default
  \texttt{7.0}. Current execution records this value in the build context and
  log. The \texttt{ProtonateWithReduce} implementation has no pH parameter, so
  this value does not alter reduce's global settings in the current code.
\item \texttt{--no-mopac}: disables MOPAC-family single-conformation
  calculators. Without this flag, MOPAC is on in this mode.
\end{itemize}

Required input:

\begin{itemize}
\item \texttt{FILE} after \texttt{--pdb}. Validation rejects an omitted path,
  a missing file, or a non-regular file.
\end{itemize}

Parse and validation diagnostics include:

\begin{lstlisting}
--pdb requires a file path: --pdb FILE
--pdb: file not found: FILE
--pdb: not a regular file: FILE
\end{lstlisting}

The reader loads the original file, calls \texttt{reduce} in BUILD mode, parses
the protonated PDB with cif++, assigns ff14SB charges, prepares force-field
charges, and computes the net charge by rounding the charge sum to the nearest
integer.

The reduce wrapper strips existing ATOM hydrogens before rebuilding them. It
adds HIS sidechain hydrogens, keeps OH/SH hydrogens, adds non-water hydrogens,
does not add water hydrogens, rotates existing OH groups, flips NQH groups, and
runs reduce optimization. The het dictionary must exist at the CMake-provided
\texttt{HET\_DICTIONARY} path or the \texttt{REDUCE\_HET\_DICT} override path.
If the input text is empty, no model is parsed, the het dictionary is missing,
or reduce returns no PDB text, this mode fails.

PDB parsing rules:

\begin{itemize}
\item \texttt{TITLE}, \texttt{REMARK}, \texttt{MODEL}, \texttt{ENDMDL}, and
  \texttt{USER} records are skipped before cif++ parsing.
\item A blank chain ID in an ATOM or HETATM record is changed to \texttt{A}.
\item An \texttt{END} record is appended if the input did not contain one.
\item Only alternate locations with an empty alt id or \texttt{A} are loaded.
\item Unknown elements are skipped.
\item Residues that do not map to a known amino-acid type are skipped.
\item If no protein atoms remain, parsing fails with
  \texttt{No protein atoms found in PDB input}.
\item Non-numeric residue sequence IDs are logged and stored as sequence
  number \texttt{0}; they are not a hard rejection in this reader.
\end{itemize}

Charge setup uses \texttt{RuntimeEnvironment::Ff14sbParams()}. If the ff14SB
flat table cannot be found, the mode fails with an ff14SB params error. If
charge resolution or charge preparation fails, the mode fails with that charge
error.

Output:

\begin{itemize}
\item \texttt{atoms\_category\_info.npy}.
\item Topology sidecars: \texttt{residues.npy}, \texttt{bonds.npy},
  \texttt{rings.npy}, \texttt{ring\_membership.npy}, and
  \texttt{extraction\_manifest.json}.
\item The identity arrays \texttt{pos.npy}, \texttt{element.npy},
  \texttt{residue\_index.npy}, and \texttt{residue\_type.npy}.
\item One set of feature arrays for the ConformationResults that attached in
  the run. APBS and AIMNet2 arrays are expected when their calculators
  succeed. MOPAC-family arrays are absent when \texttt{--no-mopac} is used.
\end{itemize}

\section{Mode 2: already-protonated PDB}

Use this mode for a PDB that already contains the intended hydrogen atoms.

\begin{lstlisting}
nmr_extract --protonated-pdb protonated.pdb --output out --aimnet2 model.jpt
\end{lstlisting}

Optional flag:

\begin{itemize}
\item \texttt{--no-mopac}: disables MOPAC-family single-conformation
  calculators. Without this flag, MOPAC is on in this mode.
\end{itemize}

Required input:

\begin{itemize}
\item \texttt{FILE} after \texttt{--protonated-pdb}. Validation rejects an
  omitted path, a missing file, or a non-regular file.
\end{itemize}

Parse and validation diagnostics include:

\begin{lstlisting}
--protonated-pdb requires a file path: --protonated-pdb FILE
--protonated-pdb: file not found: FILE
--protonated-pdb: not a regular file: FILE
\end{lstlisting}

This mode skips reduce. It parses the supplied PDB directly with the same PDB
rules as mode 1, assigns ff14SB charges, prepares force-field charges, and
writes the same output surface as mode 1.

\section{Mode 3: ORCA/AMBER-prepared pose}

Use this mode for a single tleap/AMBER-prepared pose described by a root name.

\begin{lstlisting}
nmr_extract --orca --root runs/1abc/WT --output out --aimnet2 model.jpt
\end{lstlisting}

Optional flag:

\begin{itemize}
\item \texttt{--no-mopac}: disables MOPAC-family single-conformation
  calculators. Without this flag, MOPAC is on in this mode.
\end{itemize}

Required by CLI validation:

\begin{itemize}
\item \texttt{NAME.xyz}: protonated coordinates.
\item \texttt{NAME.prmtop}: AMBER topology and charges.
\end{itemize}

The parser expands \texttt{--root NAME} exactly as:

\begin{lstlisting}
NAME.xyz
NAME.prmtop
NAME_nmr.out
\end{lstlisting}

There is no globbing and no directory search. The \texttt{NAME\_nmr.out} file
is marked optional by CLI validation, but the current runner always supplies
the expanded path to \texttt{OperationRunner}. A missing
\texttt{NAME\_nmr.out} therefore causes the run to fail when ORCA shielding is
loaded:

\begin{lstlisting}
NMR output not found: NAME_nmr.out
ORCA shielding load failed for NAME_nmr.out
\end{lstlisting}

For a first-time operator run, prepare \texttt{NAME\_nmr.out} alongside the
\texttt{.xyz} and \texttt{.prmtop}.

Parse and validation diagnostics include:

\begin{lstlisting}
--orca requires --root NAME
  root expands to: {root}.xyz, {root}.prmtop, {root}_nmr.out
--orca .xyz: file not found: NAME.xyz
--orca .prmtop: file not found: NAME.prmtop
\end{lstlisting}

ORCA/AMBER reader rules:

\begin{itemize}
\item The XYZ first line must contain a positive atom count.
\item The number of XYZ atoms must equal the number of atoms in the prmtop.
\item The prmtop must contain \texttt{ATOM\_NAME}, \texttt{RESIDUE\_LABEL},
  and \texttt{RESIDUE\_POINTER}; otherwise the reader reports
  \texttt{incomplete prmtop}.
\item Residue labels must resolve through the naming registry. Unknown
  residues are hard errors.
\item The prmtop provides the atom names, residue labels, charges, and
  protonation variants. It provides elements when \texttt{ATOMIC\_NUMBER}
  is present; otherwise the loader takes elements from the XYZ. The XYZ
  provides positions.
\item The prmtop charge source later requires \texttt{CHARGE} and
  \texttt{RADII} rows for every protein atom.
\item ORCA NMR output, when loaded, must contain the same number of nuclei as
  the conformation has atoms, and the ORCA nucleus element order must match
  the protein atom element order. A mismatch is a hard error because tensors
  would attach to the wrong atoms.
\end{itemize}

Output is the same single-conformation output surface as modes 1--2, with ORCA
shielding arrays included when the ORCA shielding result attaches.

\section{Mode 4: WT--ALA mutant pair}

Use this mode for a WT pose and its alanine mutant pose. Each side is a mode-3
ORCA/AMBER-prepared pose.

\begin{lstlisting}
nmr_extract --mutant --wt runs/1abc/WT --ala runs/1abc/ALA --output out --aimnet2 model.jpt
\end{lstlisting}

Optional flag:

\begin{itemize}
\item \texttt{--no-mopac}: disables MOPAC-family single-conformation
  calculators on both sides. Without this flag, MOPAC is on for both sides.
\end{itemize}

The parser expands each root independently:

\begin{lstlisting}
WT root:
  WT_NAME.xyz
  WT_NAME.prmtop
  WT_NAME_nmr.out

ALA root:
  ALA_NAME.xyz
  ALA_NAME.prmtop
  ALA_NAME_nmr.out
\end{lstlisting}

CLI validation requires the \texttt{.xyz} and \texttt{.prmtop} file for each
side and does not require either \texttt{\_nmr.out} file. Current execution
does pass both expanded \texttt{\_nmr.out} paths to the ORCA shielding loader.
If either file is missing or does not match its protonated geometry, the WT or
ALA side fails and no mutant comparison output is produced.

Parse and validation diagnostics include:

\begin{lstlisting}
--mutant requires --wt NAME --ala NAME
  each root expands to: {root}.xyz, {root}.prmtop, {root}_nmr.out
--mutant --wt .xyz: file not found: WT_NAME.xyz
--mutant --wt .prmtop: file not found: WT_NAME.prmtop
--mutant --ala .xyz: file not found: ALA_NAME.xyz
--mutant --ala .prmtop: file not found: ALA_NAME.prmtop
\end{lstlisting}

The runner loads both poses, runs the standard calculator sequence on both,
then computes WT--ALA mutation delta tensors when both conformations have an
ORCA shielding result. The output directory receives WT topology sidecars and
WT conformation feature arrays. The mutation delta result attaches to the WT
conformation; the runner does not write a separate full ALA output tree.

\section{Mode 5: GROMACS trajectory}

Use this mode for a GROMACS production directory with fixed file names.

\begin{lstlisting}
nmr_extract --trajectory /data/1abc/production_scan1 --stride 10 --output out --aimnet2 model.jpt
\end{lstlisting}

MOPAC is off by default in this mode. Add \texttt{--mopac} to use the FullFat
trajectory shape:

\begin{lstlisting}
nmr_extract --trajectory /data/1abc/production_scan1 --stride 10 --mopac --output out --aimnet2 model.jpt
\end{lstlisting}

To process a bounded span of raw TRR frame indices, give both window flags:

\begin{lstlisting}
nmr_extract --trajectory /data/1abc/production_scan1 --window-start 1000 --window-len 250 --stride 10 --output out --aimnet2 model.jpt
\end{lstlisting}

The mode directory must contain these exact regular files:

\begin{lstlisting}
DIR/production.tpr
DIR/production.trr
DIR/production.edr
\end{lstlisting}

The reader also looks for:

\begin{lstlisting}
DIR/../topol.top
\end{lstlisting}

\texttt{topol.top} is not required by CLI validation. If it is present and its
GROMACS-to-AMBER readback block parses, it is used to preserve residue
chemistry decisions such as CYX and HIS variants. If it is missing or cannot
be parsed, the loader falls back to the naming registry and logs that CYX/HIS
resolution may be lost.

There is no file discovery. The derivation function is
\texttt{FromProductionDir}, and it returns these paths:

\begin{lstlisting}
return {dir / "production.tpr",
        dir / "production.trr",
        dir / "production.edr"};
\end{lstlisting}

Validation checks those exact paths:

\begin{lstlisting}
e = CheckFile(files.tpr, "trajectory production.tpr"); if (!e.empty()) return e;
e = CheckFile(files.trr, "trajectory production.trr"); if (!e.empty()) return e;
e = CheckFile(files.edr, "trajectory production.edr"); if (!e.empty()) return e;
\end{lstlisting}

An \texttt{.xtc} file is not used. There is no separate explicit
\texttt{production.xtc} rejection branch; a directory containing only
\texttt{production.xtc} for coordinates fails because
\texttt{production.trr} is absent:

\begin{lstlisting}
trajectory production.trr: file not found: DIR/production.trr
\end{lstlisting}

The frame reader opens and reads \texttt{production.trr} through the TRR API:

\begin{lstlisting}
gmx_trr_open
gmx_trr_read_frame
\end{lstlisting}

It reads positions, velocities when present, and the box matrix. A TRR without
velocities is not rejected by this reader; the velocity buffer is left empty
for velocity-consuming results. A TRR frame read for processing whose atom
count differs from the TPR topology is rejected:

\begin{lstlisting}
TRR frame natoms mismatch: header=N expected=M
\end{lstlisting}

\subsection{Trajectory stride}

The current code has one trajectory cadence knob:

\begin{lstlisting}
--stride N
\end{lstlisting}

Default \texttt{N} is \texttt{1}. The parser converts \texttt{0} to
\texttt{1}. Non-numeric strings are parsed by \texttt{std::strtoull}; the code
does not currently validate that the whole token was numeric. The stride is
applied by \texttt{RunConfiguration::SetStride}. It controls frame dispatch:
with no window, frame 0 is processed, then every \texttt{N}-th TRR frame after
that is processed. With a window, dispatch starts at \texttt{window\_start}.
MOPAC, H5 time-series, per-frame NPY emission, and per-frame PDB emission all
run on exactly the dispatched frames. There is no current
\texttt{--npy-stride}, \texttt{--pdb-stride}, or \texttt{--mopac-stride}.

\subsection{Trajectory window}

Windowed trajectory extraction uses the paired flags:

\begin{lstlisting}
--window-start N --window-len M
\end{lstlisting}

Both flags must be supplied together, or neither. \texttt{N} and \texttt{M}
must be non-negative integer tokens, and \texttt{M} must be at least
\texttt{1}; invalid or partial windows fail at parse time. Omitting both
flags processes the full trajectory.

The window is a bound on the single dispatch loop over raw TRR frame indices:
\texttt{[N, N+M)}. It is not a second emission gate. \texttt{--stride} still
applies within the window, and MOPAC when enabled, H5 time-series, per-frame
NPY emission, and per-frame PDB emission all see exactly the dispatched
window-local frames. The first conformation \texttt{conf0} is the first frame
of the window; the trajectory H5 records \texttt{window\_mode},
\texttt{window\_start}, and \texttt{window\_len} attributes under
\texttt{/trajectory/}.

\subsection{Trajectory loader rules}

The TPR loader classifies leading non-water, non-ion molblocks as the protein
region. Water is a 3-atom moltype with one oxygen and two hydrogens in one
residue. An ion is a one-atom, one-residue moltype. Any non-water, non-ion
moltype after the protein/solvent region is rejected:

\begin{lstlisting}
Unsupported moltype after protein/solvent region: 'NAME' (atoms=A residues=R nmol=M). Cofactors and co-solutes are not currently supported by the trajectory loader.
\end{lstlisting}

Each protein-shaped molblock must have \texttt{nmol=1}. Replicated protein
moltypes are rejected:

\begin{lstlisting}
Protein-shape molblock 'NAME' has nmol=M (expected 1; pdb2gmx convention). Replicated protein moltypes are not supported.
\end{lstlisting}

At least one protein-shaped molblock must exist:

\begin{lstlisting}
no protein-shape molblock found in DIR/production.tpr
\end{lstlisting}

When a parsed \texttt{topol.top} readback block is used, every protein residue
must have a matching readback entry. Missing entries, out-of-range residue
indices, and mismatches between the TPR residue name and the readback comment
are hard errors. Without a readback block, residue names are resolved through
the naming registry. Unknown residue names are hard errors.

The trajectory protein is finalized from the first TRR frame. Protein positions
are converted from nm to Angstrom. Protein velocities, when present, are
converted from nm/ps to Angstrom/ps. Solvent water and ion positions are split
into the solvent side channel for solvent-aware calculators.

\subsection{Trajectory output}

Trajectory mode requires \texttt{--output DIR}. It writes:

\begin{itemize}
\item \texttt{DIR/trajectory.h5}, containing trajectory result groups,
  trajectory frame indices/times, selections, and trajectory protein atoms.
\item One top-level topology sidecar set in \texttt{DIR}.
\item Per-frame PDB files in \texttt{DIR/pdbs}. File names are
  \texttt{STEM\_fNNNNNN\_tT.T.pdb}, where \texttt{STEM} is the trajectory
  directory basename, \texttt{NNNNNN} is the raw TRR frame index padded to six
  digits, and \texttt{T.T} is the frame time in ps with one decimal place.
\item Per-frame NPY directories in \texttt{DIR/npys/frame\_NNNNNN}. Each
  directory receives the same per-conformation feature-array surface that a
  single-conformation run writes for that frame. \texttt{DIR/npys} also
  receives the per-protein sidecars on the first emitted frame.
\end{itemize}

Without \texttt{--mopac}, the run shape is
\texttt{PerFrameExtractionSet}: APBS and AIMNet2 run on every dispatched frame,
MOPAC is skipped, and vacuum Coulomb is skipped. With \texttt{--mopac}, the
run shape is \texttt{FullFatFrameExtraction}: it inherits the per-frame stack,
enables MOPAC on every dispatched frame, enables vacuum Coulomb for the
MOPAC-vs-FF14SB reconciliation probe, and adds the MOPAC-family trajectory
results.

\section{MOPAC behavior by mode}

Modes 1--4:

\begin{itemize}
\item Default: MOPAC on.
\item Meaningful switch: \texttt{--no-mopac}.
\item APBS stays on.
\item AIMNet2 stays on and requires a model.
\item Home-rolled vacuum Coulomb stays skipped in these single-conformation
  modes.
\end{itemize}

Mode 5:

\begin{itemize}
\item Default: MOPAC off.
\item Meaningful switch: \texttt{--mopac}.
\item APBS and AIMNet2 stay on.
\item \texttt{--mopac} changes the run shape from
  \texttt{PerFrameExtractionSet} to \texttt{FullFatFrameExtraction}. That also
  enables vacuum Coulomb for the FullFat reconciliation result.
\end{itemize}

\section{Out of scope and removed interfaces}

These are not supported modes or features in the current five-mode interface:

\begin{itemize}
\item propka.
\item kaml.
\item \texttt{--analysis} and \texttt{--analysis-h5}.
\item \texttt{--fleet}.
\item \texttt{--no-apbs}.
\item \texttt{--no-coulomb}.
\item \texttt{--npy-stride}, \texttt{--pdb-stride}, and
  \texttt{--mopac-stride}.
\end{itemize}

The parser's lack of a general unknown-flag rejection pass does not make these
interfaces available. If one of these flags appears beside a valid mode, it may
be ignored rather than rejected. Prepare commands using only the flags
documented above.

\section{Discrepancies found during code audit}

The guide above follows the current code behavior. The audit found these
differences between the canonical five-mode prose, CLI help/comments, and
current execution:

\begin{itemize}
\item The canonical mode block says MOPAC is an orthogonal
  \texttt{--mopac}/\texttt{--no-mopac} toggle across all five modes. Current
  parsing is not symmetric: modes 1--4 only consume \texttt{--no-mopac} and
  default on; trajectory only consumes \texttt{--mopac} and defaults off.
\item Older canonical prose said trajectory NPY and PDB emissions had
  independent strides \texttt{m} and \texttt{n}. Current code has one
  \texttt{--stride N}; processing, MOPAC when enabled, NPY emission, PDB
  emission, and H5 time-series all use that same dispatched-frame cadence.
\item The brief asks for an \texttt{.xtc}-rejected validation check. Current
  validation does not mention \texttt{.xtc}. It requires
  \texttt{production.trr}; a directory with only \texttt{production.xtc} fails
  through the missing-\texttt{production.trr} file check.
\item CLI usage text and \texttt{Validate.cpp} mark \texttt{NAME\_nmr.out} as
  optional for \texttt{--orca} and \texttt{--mutant}. Current execution always
  sets the expanded \texttt{\_nmr.out} path in \texttt{RunOptions}, so a
  missing file causes ORCA shielding load failure. The operator guidance
  therefore treats \texttt{\_nmr.out} as required for a successful mode-3 or
  mode-4 run as currently wired.
\item \texttt{--pH} exists in code and help for \texttt{--pdb}, with default
  \texttt{7.0}. It is not named in the canonical five-mode block. Current
  execution records the value but does not pass it into the reduce wrapper.
\item Common options \texttt{--output}, \texttt{--config}, \texttt{--aimnet2},
  and help flags are real CLI flags but are not part of the mode-selector list
  in the canonical five-mode block.
\end{itemize}

\end{document}
